\section{September 8th: Inclusion/Exclusion Principles, Random Variables}

\subsection{Inclusion/Exclusion Principles}

\begin{factbox}
   For any events $A$, $B$:
   \begin{align*}
      \PP(A) = \PP(A \cap B) + \PP(A \cap B^c) \\
      \PP(B) = \PP(B \cap A) + \PP(B \cap A^c) 
   \end{align*}
\end{factbox} 

Using the these identities for $\PP(A)$ and $\PP(B)$, we have the following identity:

\begin{factbox}
   For any events, $A$, $B$:
   \[ \PP(A \cup B) = \PP(A) + \PP(B) - \PP(A \cap B) \]
   Moreover, for any events $A$, $B$, $C$, we have:
   \[ \PP(A \cup B \cup C) = \PP(A) + \PP(B) + \PP(C) - \PP(A \cap B) - \PP(A \cap C) + \PP(B \cap C) - \PP(A \cap B \cap C) \]
\end{factbox}

These identities can easily be proven either using a Venn diagram or by using the identities given in Fact 3.1.1. 

And in fact, for any events $E_1, E_2, \dots, E_n$, we have the following identity:

\begin{propbox}[Inclusion/Exclusion Formula]
   For any events $E_1, E_2, \dots, E_n$:
   \[ \PP\left(\bigcup_{k = 1}^n E_k\right) = \sum_{k = 1}^{n}(-1)^{k+1}\sum_{\substack{S \subseteq [n] \\ |S| = k}}\PP\left(\bigcap_{i \in S}E_i\right) \]
\end{propbox}

\begin{proof}
   In Fact 3.1.2 we already proved the identity holds for $n = 2$ and $n = 3$. We will use induction to show that the proposition holds for all $n \geq 2$. 

   Assume the proposition holds for all $n$. For $n + 1$, we have events $E_1, \dots, E_n, E_{n+1}$. Let $A := E_1, \dots, E_n$ and $B := E_{n + 1}$. From the $n = 2$ case, we have $\PP(A \cup B) = \PP(A) + \PP(B) - \PP(A \cap B)$. By induction, we have:
   \begin{align*}
      &\PP(A) = \sum_{k = 1}^{n}(-1)^{k+1}\sum_{\substack{S \subseteq [n] \\ |S| = k}}\PP\left(\bigcap_{i \in S}E_i\right) \\
      &\PP(B) = \PP(E_{n+1}) \\
      &A \cap B = (E_1 \cup \dots \cup E_n) \cap E_{n+1} = (E_1 \cap E_{n+1}) \cup \dots \cup (E_n \cap E_{n+1}) \\
      &\PP(A \cap B) = \sum_{k = 1}^{n}(-1)^{k+1}\sum_{\substack{S \subseteq [n] \\ |S| = k}}\PP\left(\bigcap_{i \in S}E_i \cap E_{n+1}\right) 
   \end{align*}
   So using the $n = 2$ case, we have:
   \begin{align*}
      \PP(A \cup B) &= \sum_{k = 1}^{n}(-1)^{k+1}\sum_{\substack{S \subseteq [n] \\ |S| = k}}\PP\left(\bigcap_{i \in S}E_i\right) + \PP(E_{n+1}) -  \sum_{k = 1}^{n}(-1)^{k+1}\sum_{\substack{S \subseteq [n] \\ |S| = k}}\PP\left(\bigcap_{i \in S}E_i \cap E_{n+1}\right)  \\
                    &= Ask at Office Hours \qedhere
   \end{align*}
\end{proof}

\begin{example}
   Suppose there are $n$ individuals with $n$ hats. Suppose that all $n$ hats get mixed up. Each individual chooses a hat uniformly at random. What is the probability that at least one person gets their own hat back?
\end{example}
\begin{solution}
   Let $E_i$ be the event that person $i$ gets their own hat. We want to find $\PP\left(\bigcup_{i \in [n]}E_i\right)$. Using Proposition 3.1.3, we have:
   \[ \PP\left(\bigcup_{i \in [n]}E_i\right) = \sum_{k = 1}^{n}(-1)^{k+1}\sum_{\substack{S \subseteq [n] \\ |S| = k}}\PP\left(\bigcap_{i \in S}E_i\right) \]
   We know that for $n$ hats, there are $n!$ possible results from this experiment, and moreover that there are $(n-k)!$ ways to distribute the hats among $k$ individuals, hence $\PP\left(\bigcap_{i \in S}E_i\right) = (n-k)!/n!$. Also note that there are $\binom{n}{k}$ terms in the inner summation. The equation simplifies to:
   \begin{align*}
      \PP\left(\bigcup_{i \in [n]}E_i\right) &= \sum_{k = 1}^{n}(-1)^{k+1}\binom{n}{k}\frac{(n-k)!}{n!} \\
                                             &= \sum_{k = 1}^{n}(-1)^{k+1}\frac{1}{k!} \\
                                             &= 1 - \sum_{k = 0}^{n}\frac{(-1)^k}{k!} \\
                                             &\approx 1 - \frac{1}{e} \qedhere
   \end{align*}
\end{solution}

\subsection{Random Variables}




